\chapter{One Circuit, Four Views}
\label{ch:fourviews}
Part~II built the generic machinery (systems seen as differential equations,
state-space models, transfer functions, and frequency responses), and the
previous two chapters supplied the physical laws and units. This short chapter
states the organizing pattern the rest of Part~III follows. Its claim is simple:
\emph{every circuit in this book is one of those systems}, and the way to
understand a circuit is to look at it through the same set of views until they
agree.

\section{The Four Views}
A circuit can be read as:
\begin{enumerate}
\item a \textbf{physical network} governed by conservation laws (KCL, KVL) and
element laws;
\item a \textbf{differential equation} in the natural energy-storage variables;
\item a \textbf{state-space model} whose eigenvalues are the circuit's poles;
\item a \textbf{phasor / frequency-response model}, the transfer function on the
\(\jj\omega\) axis.
\end{enumerate}
A fifth lens---\textbf{energy} storage and dissipation---runs alongside as a
consistency check. All of these must agree; if they do not, a sign convention, a
reference direction, or an assumption needs attention.

\section{The Recipe for Part III}
The next chapters each apply the same recipe to a specific circuit, and it is
worth stating once so the pattern is visible:
\begin{enumerate}
\item draw the schematic and name the state variables (capacitor voltages,
inductor currents);
\item apply KVL/KCL and the element laws to get the differential equation;
\item write the state-space model and read off the poles;
\item \textbf{match it to the generic first- or second-order system of
\cref{ch:linsys}}---identify \(\tau\), or \(\omega_0\) and \(\zeta\)---so all of
Part~II's results apply at once;
\item interpret the frequency response (Bode) and the pole motion (root locus);
\item close with the energy balance and a worked example.
\end{enumerate}

\section{The Map from Circuits to Systems}
Three circuits carry most of the exam's analog content, and each is a generic
system in disguise:

\begin{center}
\begin{tabularx}{\textwidth}{@{}L{0.15\textwidth}L{0.19\textwidth}L{0.24\textwidth}Y@{}}
\toprule
Circuit & Order / state(s) & Generic system & Key parameters \\
\midrule
Series RC & first / \(v_C\) & \(\tau\dot y+y=Ku\) & \(\tau=RC\), pole \(-1/RC\) \\
Series RL & first / \(i_L\) & \(\tau\dot y+y=Ku\) & \(\tau=L/R\), pole \(-R/L\) \\
Series RLC & second / \(v_C,i_L\) & \(\ddot y+2\zeta\omega_0\dot y+\omega_0^2 y=\omega_0^2u\) & \(\omega_0=\tfrac{1}{\sqrt{LC}}\), \(\zeta=\tfrac{R}{2}\sqrt{C/L}\) \\
\bottomrule
\end{tabularx}
\end{center}

Once a circuit is matched to its row, everything from \cref{ch:linsys}---step
response, time constant or damping ratio, pole geometry, resonant peak, root
locus---transfers directly. The circuit chapters that follow do exactly this,
one circuit at a time.
